\centerline{\bf \Huge Тренировка региона Эйлера}

\begin{flushright}
        \scriptsize{}
\end{flushright}

\problem{8.1}
Операция удвоения цифры натурального числа состоит в умножении этой цифры на 2 (если это произведение оказывается двузначным, то цифра в следующем разряде числа увеличивается на 1 , как при сложении «в столбик»). Например, из числа 9817 удвоениями цифр 7, 1, 8 и 9 можно получить числа 9824, 9827, 10617 и 18817 соответственно. Можно ли из числа $22 . .22$ (20 двоек) несколькими такими операциями получить число 22... 22 (21 двойка)?


\problem{8.2}
В пещере собрались 100 гномов - по 10 гномов из 10 разных кланов. Каждый из них - рыцарь, всегда говорящий правду, или лжец, который всегда лжет. Каждый из собравшихся назвал клан, из которого, по его мнению, на собрание пришли одни лжецы. Оказалось, что каждый из 10 кланов назвало ровно 10 гномов. Докажите, что лжецов в пещере не меньше, чем рыцарей.

\problem{8.3}
Числа $x, y, z$ таковы, что $x>y^2+z^2, y>z^2+x^2, z>x^2+y^2$. Докажите, что каждое из чисел $x, y, z$ меньше $\frac{1}{2}$.

\problem{8.4}
В каждой клетке доски $2 \times 200$ лежит по рублёвой монете. Даша и Соня играют, делая ходы по очереди, начинает Даша. За один ход можно выбрать любую монету и передвинуть её: Даша двигает монету на соседнюю по диагонали клетку, Соня - на соседнюю по стороне. Если две монеты оказываются в одной клетке, одна из них тут же снимается с доски и достаётся Соне. Соня может остановить игру в любой момент и забрать все полученные деньги. Какой наибольший выигрыш она может получить, как бы ни играла Даша?

\problem{8.5}
На биссектрисе угла $A B C$ отмечена точка $D$. На отрезке $A B$ отмечена точка $E$, а на отрезке $B C$ - точка $F$, причём $A B=D E$ и $B C=D F$. Докажите, что из отрезков $A D, C D$ и $E F$ можно сложить треугольник.